\section{Related Work}
\label{sec:related_work}

\subsection{Looped Transformers}

We use the term \emph{looped model} to denote an architecture that reuses a learned internal operator within a single forward computation, rather than repeatedly calling an otherwise complete model. This design lineage predates contemporary language models: Neural GPUs repeatedly apply a shared convolutional transition to learn algorithms, Adaptive Computation Time (ACT) learns how many recurrent updates to execute, Universal Transformers (UTs) tie self-attention and transition blocks across depth, and Deep Equilibrium Models (DEQs) avoid explicit unrolling by solving for a fixed point of a weight-tied transformation \citep{kaiser2016neural,graves2016adaptive,dehghani2018universal,bai2019deep,bai2020multiscale}. Work on training implicit models and on generalist neural algorithmic learners further established practical optimization methods and reusable processors for effectively unbounded or task-shared computation. These are direct architectural antecedents of looped Transformers, not merely generic recurrent or test-time-compute methods.

Across Transformer families, layer tying, recurrent state, adaptive depth, and memory reuse also appear in ALBERT, tied Transformers, depth-adaptive Transformers, feedback memory, recurrent-memory Transformers, and block-recurrent Transformers \citep{lan2019albert,xia2019tied,elbayad2020depth,fan2020feedback,bulatov2022rmt,hutchins2022blockrecurrent}. Modern variants make the shared-depth computation itself the central scaling axis. MoEUT recurrently reuses fine-grained expert groups, Relaxed Recursive Transformers supplement a repeated block with depth-specific low-rank adapters, and implicit state-space language models iterate a shared transition toward a fixed point while preserving substantial training parallelism \citep{csordas2024moeut}. Thus, rather than relying on ordinary depth scaling, looped Transformers learn a reusable transition $h_{t+1}=F_{\theta}(h_t,x,e_t)$, where $e_t$ may encode the loop index, a halting state, or an injected copy of the input. Effective inference depth can then vary without introducing an independently parameterized layer at every step.

\subsection{Inductive Biases}

Looped depth imposes several biases that are weak or absent in a stack of independent Transformer layers. \emph{First, recurrent operator sharing} encourages a reusable update rule rather than a sequence of layer-specific feature maps. This is the defining recursive bias of UTs and tied-depth models, and theoretical analyses connect it to fixed-point iteration, gradient-based in-context learning, normalized gradient descent, and power iteration \citep{dehghani2018universal,takase2023lessons,yang2023looped,gatmiry2024can,chen2025bypassing,wu2026powermethod}. \emph{Second, shared depth decouples effective computation from parameter count}: a compact recurrent core can be unrolled for more steps, favoring iterative refinement over memorizing a separate computation in each layer \citep{lan2019albert,saunshi2025latentthoughts,geiping2025scalinglatent,zhu2025scalinglatent,schwethelm2026much}.

\emph{Third, looping induces an iterative-algorithm bias}. When a target computation is naturally expressed by repeated application of a small rule set, looped models can emulate programs, optimization procedures, graph algorithms, and structured algorithmic processors, often improving length or depth extrapolation when the loop count adapts to problem size \citep{giannou2023looped,backdeluca2024simulation,gao2025algoformer,fan2024looped,xu2025cotloop}. \emph{Fourth, weight sharing creates a knowledge re-access and compositionality bias}: after each new bridge entity is inferred, the same retrieval, binding, and update machinery can be applied again. This directly targets failures of implicit multi-hop composition in ordinary Transformers and motivates architectures that carry discrete and continuous states or align latent and explicit reasoning \citep{wang2024grokking,biran2024hopping,guo2025twohop,yao2025implicit,kohli2026loopthinkgeneralize,fu2026discoloop,fan2026lotus}. \emph{Fifth, recurrence provides a latent scratchpad and an adaptive-computation bias}. Additional computation occurs in the hidden state rather than through the generation of extra verbal tokens, while timestep conditioning, token-wise routing, dynamic halting, fixed-point convergence, and shortcut consistency allocate more computation to harder inputs \citep{saunshi2025latentthoughts,geiping2025scalinglatent,bay2025mixture,fu2025thinkathard,jeddi2026loopformer,movahedi2026fixedpoint}. These benefits are not automatic: looped models can drift, overthink, or collapse to shallow computation unless their state injection, normalization, supervision, and stopping rules are designed for stable iteration \citep{chowdhury2024investigating,labovich2026stability,rauba2026tarm,moosa2026dynamiccompute}.

\subsection{Theory, Mechanisms, and Scaling}

Formal results characterize both the capabilities and the limits of shared-depth computation. Looped Transformers can execute instruction-level programs, simulate latent chains of thought, and achieve better approximation rates through timestep modulation \citep{giannou2023looped,saunshi2025latentthoughts,xu2025expressive}. More specialized analyses show that a recurrently reused attention layer can learn normalized-gradient updates for in-context logistic regression or implement the power method under layer normalization \citep{wu2026powermethod}. At the same time, compressed recurrent states create a memory-budget separation relative to full sequence-state chain-of-thought methods, and set-valued or graph computations require careful distinctions among fixed points, convergence, and halting. A complementary reinforcement-learning analysis formalizes how additional internal computation changes the class of compute-bounded policies.

A second line of work studies optimization, stability, and scaling with loop count. Fixed-point analyses identify conditions under which recurrent states remain reachable, input-dependent, and trainable; residual-scaling theory argues that correlated tied-block updates require loop-aware scaling; and iso-depth studies quantify the exchange rate among recurrent passes, unique depth, and training compute \citep{labovich2026stability,wang2026residualscaling,schwethelm2026much}. Two-scale latent-dynamics analyses derive convergence-sensitive early-exit criteria, while controlled studies of adaptive computation show that learned loop allocation can correlate with token difficulty without necessarily extrapolating to longer inputs \citep{pappone2025twoscale,moosa2026dynamiccompute}. Memory tokens and ACT initialization can determine whether a UT enters a non-trivial reasoning regime, and hierarchical and flat recurrence can exhibit materially different optimization behavior even under shared-weight controls \citep{sapunov2026utm,han2026hierarchicalflat}. Compression studies likewise show that preserving per-cell accuracy or local states is insufficient if quantization or pruning damages the recursive trajectory needed for exact solutions.

Mechanistic work examines what is represented across recurrent steps. Probing studies find that recurrent language models often approach loop-specific fixed points or replay feed-forward-like inference stages, while evidence for a literal latent chain-of-thought is mixed and depth-dependent \citep{blayney2026mechanistic,lu2025latentcot}. Shared recurrent modules can nevertheless specialize to perform distinct functional roles through asymmetric state identities, and interaction-locality measurements reveal how repeated local updates accumulate into global puzzle structure. Other probes recover relational preference information from differences between loop states, and analyses of tabular foundation models use the observed layer redundancy to motivate a single repeatedly applied layer. Finally, recurrent-depth models can unlock systematic multi-hop generalization, but autoregressive controls show that the benefit depends on where compute is placed rather than on recursion alone \citep{kohli2026loopthinkgeneralize,rauba2026tarm}.

\subsection{Architectures, Training Objectives, and Inference}

Recent looped language models treat recurrence as a third scaling axis alongside parameter count and token count. Huginn and Ouro pretrain recurrent-depth language models whose shared cores can be unrolled for variable test-time depth, while retrofitted recurrence and LoopUS convert pretrained feed-forward models into latent-refinement systems without training a recurrent model entirely from scratch \citep{geiping2025scalinglatent,zhu2025scalinglatent,mcleish2025retrofitted,park2026loopus}. HRM-Text and CHERRY explore recurrent compression of deep language models, Hyperloop repeats only a middle block with cross-loop connections, and LoopMoE, sparse looped layers, and universal expert pools combine iterative depth with conditional capacity \citep{wang2026hrmtext,kwon2026cherry,zeitoun2026hyperloop,chen2026loopmoe,lee2026sparse}. MoEUT and Relaxed Recursive Transformers provide earlier expert-routing and low-rank mechanisms for recovering expressivity under parameter sharing \citep{csordas2024moeut}.

Dynamic-depth architectures decide not only \emph{how} to update a state but also \emph{where} and \emph{for how long} to recur. Mixture-of-Recursions routes tokens to different depths; CoTFormer exposes earlier recurrent representations and learns a compute-budgeted router; AdaPonderLM and Think-at-Hard allocate extra latent iterations selectively; and Chain-of-Layers methods skip or repeat pretrained layers at test time \citep{bay2025mixture,mohtashami2025cotformer,fu2025thinkathard,li2026programlayers}. LoopFormer regularizes trajectories across sampled depths, adaptive-loop models couple halting with external memory, and ChainGPT, MoDr, and depth-recurrent attention mixtures enrich the recurrent state transition through multi-rank updates, branch routing, or mixtures of sequence and depth attention \citep{jeddi2026loopformer,frey2026adaptive,zheng2026chaingpt,zhang2026modr}. SpiralFormer instead changes the resolution schedule across repeated applications, and subgoal-persistence models study when a hierarchical reasoner should re-plan rather than update the same latent plan at every step \citep{yu2026spiralformer,chadha2026replan}.

A parallel line of work targets stable and efficient unrolling. Parallel Loop Transformers share representations across loops, and LT2 replaces quadratic attention with a linear-time looped design \citep{wu2025parallellooptransformerefficient,deng2026lt2}. MELT keeps cache memory independent of loop depth; LASER compresses recursive activations during training; and Hyperloop and CHERRY reduce the number of distinct parameter blocks \citep{cakar2026laser,zeitoun2026hyperloop,kwon2026cherry}. Parcae constrains recurrent dynamics and derives scaling laws, CART anchors each iteration to context, fully looped signal routing stabilizes high loop counts, stochastic stopping improves extrapolation across unseen depths, and stability-aware recurrent training reduces drift under test-time scaling \citep{prairie2026parcae,capps2026cart,fu2026fullylooped,kuo2026stochasticstopping}. Training-free looping explores whether a pretrained model can be recurrently reused without architectural retraining \citep{chen2026trainingfree}.

Several methods supervise the \emph{trajectory} rather than the final answer alone. LoopRPT applies reinforcement pre-training across latent iterations, RLTT distributes reward over the latent trajectory, denoising recursion trains models to repeatedly correct corrupted targets, and Generative Recursive Reasoning expands recurrence into stochastic multi-trajectory generation \citep{tang2026looprpt,cameron2026denoisingrecursion,baek2026gram}. Probabilistic TRM injects noise at inference and selects among recursive trajectories, whereas LoopFormer uses shortcut consistency to align different computation budgets \citep{sghaier2026ptrm,jeddi2026loopformer}. Looped diffusion language models, fixed-point masked generative models, equilibrium reasoners, and attractor models connect explicit recurrence to denoising or convergence-based generation \citep{miele2026fixedpointmasked,feinashley2026attractor}. These approaches make intermediate-state quality, convergence geometry, and trajectory diversity first-class training targets rather than incidental by-products of depth.

\subsection{In-Context, Algorithmic, and Compositional Generalization}

In-context learning can be interpreted as executing an implicit learning algorithm over demonstrations. Bayesian, gradient-descent, preconditioned-gradient, linear-model, and causal-structure accounts motivate studying whether each recurrent step implements one additional internal update \citep{xie2022explanation,akyurek2023learning,vonoswald2023transformers,ahn2023transformers,li2023algorithms,zhang2023trained,nichani2024causal}. Looped Transformers make this correspondence explicit: they learn data-fitting algorithms with substantially fewer unique parameters, implement multi-step gradient procedures, and can use distinct preprocessing, looping, and postprocessing stages \citep{yang2023looped,gatmiry2024can,chen2025bypassing,gao2025algoformer}. Program-simulation results and normalized-gradient analyses further connect shared depth to reusable algorithmic primitives \citep{giannou2023looped}.

The same prior is relevant to length and compositional extrapolation. Neural GPUs and generalist recurrent processors learn repeated algorithmic updates, while looped Transformers simulate graph algorithms and obtain strong length generalization when the number of recurrent steps grows with the instance \citep{kaiser2016neural,backdeluca2024simulation,fan2024looped}. Timestep modulation, multi-resolution recursion, and depth-recurrent compositional models address the failure of a single stationary update to remain useful far beyond the trained depth \citep{xu2025expressive,yu2026spiralformer,chen2026thinking}. DiscoLoop carries both discrete embeddings and continuous hidden states across hops, and LOTUS aligns recurrent latent blocks with explicit chain-of-thought computation \citep{fu2026discoloop,fan2026lotus}. These results complement evidence that ordinary Transformers often memorize local facts or short computations yet fail to compose them systematically out of distribution \citep{biran2024hopping,yao2025implicit,dziri2023faith}.

\subsection{Abstract Reasoning}

ARC-style tasks are a natural stress test because they require inducing a new transformation from a few demonstrations rather than recalling a fixed skill \citep{chollet2019measure,chollet2025arcagi2}. HRM obtains substantial effective depth through coupled high- and low-level recurrent modules, TRM reduces the design to a tiny repeatedly applied network, and the URM study attributes much of the gain to Universal-Transformer-style recurrence and nonlinear depth computation \citep{wang2025hierarchical,jolicoeurmartineau2025trm,gao2025universal}. Controlled analyses caution that hierarchy, augmentation, identity conditioning, majority voting, and competition-time adaptation can materially affect reported scores \citep{ge2025hrmperspectives,ren2026reasoningguessing,royeazar2025trm,mcgovern2025testtime}.

Follow-up designs explore complementary ways to improve recursive abstract reasoning. CosmicFish-HRM adapts hierarchical recurrence to compact language models; Recursive Inference Machines generalize generator--solver recursion; and Fixed-Point Reasoners halt upon convergence rather than at a preset depth \citep{lakkapragada2026cosmicfish,komisarczyk2026rim,movahedi2026fixedpoint}. Generative, probabilistic, and denoising recursive models introduce trajectory diversity or iterative corruption-and-repair curricula \citep{baek2026gram,sghaier2026ptrm,cameron2026denoisingrecursion}. Equilibrium and attractor reasoners learn convergent solution dynamics, whereas Tiny Autoregressive Recursive Models test whether comparable compute is better spent on recurrence or ordinary autoregressive depth \citep{feinashley2026attractor,rauba2026tarm}. Memory-augmented UTs, LoopViT, and interaction-locality analyses expose complementary depth--state, visual-recurrence, and mechanism-level views \citep{sapunov2026utm,shu2026loopvit}. Overall, this literature supports recurrence as an architectural prior for iterative abstraction, while showing that successful test-time scaling depends on stable state transitions and informative intermediate supervision rather than loop count alone \citep{chollet2026arcprize,vahdati2026arcprogress}.
